\clearpage
\thispagestyle{empty}
\pdfbookmark[0]{Magnetnadel in Kopenhagen und Mailand. Ausgabe und Quellen}{copenhagen-milan-edition}
\begin{center}\Large\textbf{Magnetnadel in Kopenhagen und Mailand\\5. und 6. November 1834}\end{center}

Carl Friedrich Gauß’ vollständiger Text aus den \textit{Werken}, Band V (Göttingen, 1867), S.~537--540, liegt hier als KI-Übertragung in modernes Deutsch mit KI-erstelltem Satz vor. Enthalten ist auch Gauß’ Fußnote zu Christopher Hansteens Karten der magnetischen Intensität.

Der Erstdruck erschien in Heinrich Christian Schumachers \textit{Astronomischen Nachrichten}, Band~12, Nr.~276, 21. März 1835, Sp.~185--188. Ihm ist die im Text erwähnte Kupfertafel beigegeben. Die obere Kurve zeigt Kopenhagen, die untere Mailand, jeweils mit eigener Skala.

\medskip\noindent
\href{https://gdz.sub.uni-goettingen.de/id/PPN236006339}{Deutsche Werkausgabe (1867)}\par
\noindent\href{https://www.digitale-sammlungen.de/view/bsb10049437?page=103}{Erstdruck (1835)}\par
\noindent\href{https://www.digitale-sammlungen.de/view/bsb10049437?page=231}{Originale Kupfertafel (Bayerische Staatsbibliothek)}
\clearpage
